%!TEX root = ../template.tex
%%%%%%%%%%%%%%%%%%%%%%%%%%%%%%%%%%%%%%%%%%%%%%%%%%%%%%%%%%%%%%%%%%%%
%% abstract-en.tex
%% NOVA thesis document file
%%
%% Abstract in English([^%]*)
%%%%%%%%%%%%%%%%%%%%%%%%%%%%%%%%%%%%%%%%%%%%%%%%%%%%%%%%%%%%%%%%%%%%

\typeout{NT FILE abstract-en.tex}%

In software development, requirements engineering plays a critical role in ensuring
that software products meet stakeholder expectations and standards.
However, the manual writing and thinking of requirements gathering, refinement,
 and validation introduces the risk of misunderstanding, incomplete requirements,
  and communication breakdowns between stakeholders and development teams.

% 

Although AI techniques have been explored in various stages of software development,
 their application in automating and improving requirements engineering is still a new 
 field of study. This thesis aims to explore the use of Generative AI (GAI) to automate the identification, clarification, and specification of requirements, ultimately improving the efficiency and accuracy of the RE process.




Concerning its contents, the abstracts should not exceed one page and may answer the following questions (it is essential to adapt to the usual practices of your scientific area):

\begin{enumerate}
  \item What is the problem?
  \item Why is this problem interesting/challenging?
  \item What is the proposed approach/solution/contribution?
  \item What results (implications/consequences) from the solution?
\end{enumerate}

% Palavras-chave do resumo em Inglês
% \begin{keywords}
% Keyword 1, Keyword 2, Keyword 3, Keyword 4, Keyword 5, Keyword 6, Keyword 7, Keyword 8, Keyword 9
% \end{keywords}
\keywords{
  One keyword \and
  Another keyword \and
  Yet another keyword \and
  One keyword more \and
  The last keyword
}
